% from ega2-1-fr.tex
\oldpage[II]{180}
$\tau$ de (\hyperref[II.8.8.1.2-fr]{8.8.1.2})~: il y a par suite un entier
$d>0$, un sous-$\mathcal{O}_Y$-Module quasi-cohérent de type fini
$\mathcal{E}$ de $\mathcal{S}_d$ tel que si $\mathcal{S}'$ est la
sous-$\mathcal{O}_Y$-Algèbre de $\mathcal{S}$ engendrée par $\mathcal{E}$ et
$\tau'=\tau\circ f^*(\varphi)$ ($\varphi$ injection canonique
$\mathcal{S}'\to\mathcal{S}$), $r'=r_{\mathcal{L},\tau'}$ soit une immersion
\[
  X\longrightarrow P'=\operatorname{Proj}(\mathcal{S}').
\]
En outre, comme $\varphi$ est injectif, $r'$ est aussi une \emph{immersion
dominante} (\hyperref[II.3.7.6-fr]{3.7.6})~; le même raisonnement que pour $r$
montre donc que $r'$ est une \emph{immersion fermée surjective}~; comme $r'$ se
factorise en
$X\xrightarrow{r}\operatorname{Proj}(\mathcal{S})
\xrightarrow{\Phi}\operatorname{Proj}(\mathcal{S}')$, où
$\Phi=\operatorname{Proj}(\varphi)$, on en conclut que $\Phi$ est aussi une
\emph{immersion fermée surjective}. Mais cela entraîne que $\Phi$ est un
\emph{isomorphisme}~; en effet, on peut se ramener au cas où
$Y=\operatorname{Spec}(A)$ est affine, $\mathcal{S}=\widetilde{S}$,
$\mathcal{S}'=\widetilde{S'}$, où $S$ est une $A$-algèbre graduée, $S'$ une
sous-algèbre graduée de $S$. Pour tout élément homogène $t\in S'$,
$S'_{(t)}$ est alors un sous-anneau de $S_{(t)}$~; si on revient à la
définition de $\operatorname{Proj}(\varphi)$
(\hyperref[II.2.8.1-fr]{2.8.1}), on voit qu'on est ramené à prouver que si
$B'$ est un sous-anneau d'un anneau $B$, et si le morphisme
$\operatorname{Spec}(B)\to\operatorname{Spec}(B')$ correspondant à l'injection
canonique $B'\to B$ est une immersion fermée, ce morphisme est nécessairement
un \emph{isomorphisme}~; mais cela résulte de
(\hyperref[I.4.2.3-fr]{I, 4.2.3}).

On a en outre
$\Phi^*(\mathcal{O}_{P'}(n))=\mathcal{O}_P(n)$
(\hyperref[II.3.5.2-fr]{3.5.2, (ii)} et
\hyperref[II.3.5.4-fr]{3.5.4}), donc
$r'^*(\mathcal{O}_{P'}(n))$ est isomorphe à $\mathcal{L}^{\otimes n}$
(\hyperref[II.4.6.3-fr]{4.6.3}). Posons
$\mathcal{S}''=\mathcal{S}'^{(d)}$, de sorte
(\hyperref[II.3.1.8-fr]{3.1.8, (i)}) que $X$ s'identifie canoniquement à
$P''=\operatorname{Proj}(\mathcal{S}'')$ et
$\mathcal{L}''=\mathcal{L}^{\otimes d}$ à $\mathcal{O}_{P''}(1)$
(\hyperref[II.3.2.9-fr]{3.2.9, (ii)}).

Cela étant, si $C''=\operatorname{Spec}(\mathcal{S}'')$,
$\mathcal{S}_{P''}^{\geqslant}=
\bigoplus_{n\geqslant 0}\mathcal{O}_{P''}(n)$ s'identifie à
$\bigoplus_{n\geqslant 0}\mathcal{L}''^{\otimes n}$, donc
$C_{P''}=\operatorname{Spec}(\mathcal{S}_{P''}^{\geqslant})$ à
$\mathbf{V}(\mathcal{L}'')$~; on sait d'autre part
(\hyperref[II.8.7.3-fr]{8.7.3}) que $C_{P''}$ est $C''$-isomorphe à
$\operatorname{Proj}(\mathcal{S}''^{\natural})$~; par définition de
$\mathcal{S}''$, $\mathcal{S}''^{\natural}$ est engendré par
$\mathcal{S}_1''^{\natural}$ et $\mathcal{S}_1''^{\natural}$ est de type fini
sur $\mathcal{S}_0''^{\natural}=\mathcal{S}''$
(\hyperref[II.8.2.10-fr]{8.2.10, (i) et (iii)}), donc
$\operatorname{Proj}(\mathcal{S}''^{\natural})$ est \emph{projectif} sur
$C''$ (\hyperref[II.5.5.1-fr]{5.5.1}). Considérons alors le diagramme
\begin{equation}
\label{II.8.8.4.1-fr}
  \xymatrix{
    \mathbf{V}(\mathcal{L})\ar[r]^{g}\ar[d]_{u} &
    \operatorname{Spec}(\mathcal{S})=C\ar[d]^{v}\\
    \mathbf{V}(\mathcal{L}'')\ar[r]_{g''} &
    \operatorname{Spec}(\mathcal{S}'')=C''
  }
\tag{8.8.4.1}
\end{equation}
où $g$ et $g''$ correspondent par (\hyperref[II.1.5.6-fr]{1.5.6}) aux
$f$-morphismes canoniques
\[
  \mathcal{S}\longrightarrow
  \bigoplus_{n\geqslant 0}\mathcal{L}^{\otimes n}
  \quad\text{et}\quad
  \mathcal{S}''\longrightarrow
  \bigoplus_{n\geqslant 0}\mathcal{L}''^{\otimes n}
\]
(\hyperref[II.3.3.2.3-fr]{3.3.2.3}) (voir
(\hyperref[II.8.8.5-fr]{8.8.5}) ci-dessous), $v$ au morphisme d'inclusion
$\mathcal{S}''\to\mathcal{S}$, et $u$ au morphisme d'inclusion
$\bigoplus_{n\geqslant 0}\mathcal{L}^{\otimes nd}\to
\bigoplus_{n\geqslant 0}\mathcal{L}^{\otimes n}$~; il est immédiat
(\hyperref[II.3.3.2-fr]{3.3.2}) que ce diagramme est commutatif. Nous venons
de voir que $g''$ est un morphisme projectif~; d'autre part, $u$ est un
morphisme \emph{fini}. En effet, la question étant locale sur $X$, on peut
supposer que $X$ est affine d'anneau $A$ et que $\mathcal{L}=\mathcal{O}_X$~;
tout revient alors à remarquer que l'anneau $A[T]$ est un module de type fini
sur son sous-anneau $A[T^d]$ ($T$ indéterminée). Comme $Y$ est un schéma
quasi-compact, et que $C''$ est affine sur $Y$, $C''$ est aussi un schéma
quasi-compact,

